\honest{Configurations and Amplitude}{Eliezer Yudkowsky}{2008}

If the world is not made of billiard balls or waves, what is it made of? Take a photon, a half-silvered mirror and two detectors. Half the time one detector clicks, half the time the other. The early scientists thought the mirror reflected the photon half the time. ``Ha, ha!'' I say. \nb{For this experiment, that reading predicts exactly what is seen. It fails only in the next experiment.}

Imagine a program that is the experiment. It holds configurations, such as ``a photon heading toward A,'' and each stores a complex number, an amplitude. The mirror passes the amplitude on multiplied by 1 when the photon goes straight and by $i$ when it turns. \nb{The standard symmetric beam splitter also divides both by $\sqrt{2}$. Without that, the totals double at each mirror; the post's results survive because it compares only paths through the same number of mirrors.} We cannot see amplitudes. A ``magical measuring tool,'' counting clicks, gives the ratios of their squared moduli. Why this is so, I do not consider ``yet.'' \nb{In ``Many Worlds, One Best Guess'' Yudkowsky later wrote: ``There is no known reason for the Born probabilities.''}

Now two half-mirrors and two full mirrors. I work out the amplitudes: at E, $i$ and $-i$ cancel; at F, $1+1=2$. No photon reaches E. \nb{The arithmetic is right, and so is the physics: this is the Mach--Zehnder interferometer.} A ball deflected at random would reach each detector half the time. Perhaps the mirror marks the photon so that it turns the other way next time? Occam's razor forbids it, and so does a third experiment: block one path, and both detectors click equally again. \nb{Of the photons that arrive. Half of all photons now stop at the block.} This ``cannot possibly'' happen with a little billiard ball. \nb{It can in pilot-wave theory, where the particle takes one path but is guided by a wave on both; that theory makes the same predictions.}

Probabilities are never negative, so they cannot cancel. Amplitudes are not degrees of belief, and configurations are not possible worlds. What could have happened but did not affects nothing. ``What affects the world is real,'' and amplitudes have visible effects, so they are real. \nb{That amplitudes are causes, and not a tool for predicting what agents will see, is the disputed premise. QBists do not say amplitudes are degrees of belief; they say the quantum state as a whole specifies an agent's credences about outcomes.}

Real configurations, I add, cover ``all the particles in the universe,'' in a continuous space. \nb{That is quantum mechanics with a fixed number of particles. Photons belong to quantum field theory.}
